\chapter{Analysis}
\label{ch:analysis}

This chapter establishes the foundation for all subsequent design and implementation decisions.
It begins by defining the problem and the target audience whose needs motivated this project, then
describes the intended architecture at a conceptual level before any technology is named. The formal
functional and non-functional requirements are stated next, followed by a use-case view grouped by
actor. The chapter closes by examining existing platforms and identifying the integration gap that
none of them currently fills.

% ----------------------------------------------------------------
\section{Problem Statement and Target Audience}
\label{sec:problem-statement}

The primary target audience of StudySpace consists of two groups whose workflows are closely
coupled: \emph{students} and \emph{instructors} at a higher-education faculty. Students are the
consumers and potential contributors of course content. Their typical lifecycle within a single
course involves browsing uploaded materials, asking questions about specific documents, annotating
or extending study notes independently, and occasionally contributing improvements that could
benefit their peers. Instructors are the publishers and curators of that content. Their lifecycle
involves creating a course structure, uploading materials into named sections, managing which
students are enrolled, and reviewing any student contributions before deciding whether to make them
visible to the whole class.

The Faculty of Information Technology at CTU in Prague serves as the representative context. A
student's daily workflow involves switching between multiple systems: the faculty course portal or
Moodle for lecture slides and reading assignments, Microsoft Teams or SharePoint for live sessions
and recordings, Progtest or GitLab for programming assignments and grading, and Discord, WhatsApp,
or email for informal peer discussion. This constant context-switching leads to scattered
information, loss of peer-generated knowledge, and extra effort spent locating the right tool for a
given task. Discussions about a specific lecture slide happen in a general chat channel with no link
to the actual material. Student notes and annotations remain in personal files and never contribute
back to the course. There is no structured way for a student to propose improvements to lecture
materials, and AI tools are used in isolation without any connection to the course content.

StudySpace aims to unify course management, student collaboration, contextual discussion, and AI
assistance into a single platform, so that the interactions between these activities are preserved
and visible to all participants.

% ----------------------------------------------------------------
\section{Architecture Overview}
\label{sec:architecture-overview}

\subsection{System Architecture}
\label{subsec:system-architecture}

At the highest level, StudySpace is a web application that separates the browser-based user
interface from the server-side application logic and the persistent data store. The user interface
runs entirely in the browser and communicates with the backend over two channels: a standard
request-response channel for operations such as loading course content, uploading materials, and
submitting contribution proposals, and a persistent real-time channel for events that must reach
all connected participants without a manual page refresh, such as chat messages and study-session
activity feeds.

The backend is a single application that handles all business logic, enforces access control,
processes file uploads, calls the external AI API, and broadcasts real-time events. The data store
holds all persistent state, including user accounts, course structures, Private Workspace records,
Merge Proposals, and conversation history.

All three tiers run as isolated containers managed by a single container orchestration file, so the
complete application can be started with one command and requires no manual environment
configuration. The final architecture diagram is presented in section~\ref{sec:system-architecture-final}
after the technology choices have been made.

% TODO: Insert high-level system architecture diagram here
% \begin{figure}[htbp]
%     \centering
%     \includegraphics[width=0.9\textwidth]{images/fig_2_1_system-architecture.png}
%     \caption{Conceptual three-tier system architecture of StudySpace}
%     \label{fig:conceptual-architecture}
% \end{figure}

\subsection{Logical Architecture}
\label{subsec:logical-architecture}

The server-side application follows the Model-View-Controller (MVC) layered pattern. The
\emph{controller layer} receives HTTP requests, delegates to the appropriate service, and serialises
the response. The \emph{service layer} contains all business logic, orchestrates data access, and
enforces domain rules such as ownership checks and status transitions. The \emph{repository layer}
provides data access through named query methods and custom queries for complex lookups. Entity
classes map the relational schema to Java objects. Data Transfer Objects (DTOs) define the API
contract and decouple the internal entity structure from the shape of the JSON responses, so the
two can evolve independently.

% TODO: Insert MVC layered architecture diagram here
% \begin{figure}[htbp]
%     \centering
%     \includegraphics[width=0.85\textwidth]{images/fig_2_2_logical-architecture.png}
%     \caption{Logical (MVC) architecture of the StudySpace backend}
%     \label{fig:logical-architecture}
% \end{figure}

% ----------------------------------------------------------------
\section{Functional Requirements}
\label{sec:requirements}

The following requirements are derived directly from the problem statement and the target audience
description in section~\ref{sec:problem-statement}. Each requirement is identified by an ID that
is referenced throughout the remainder of the thesis.

\bigskip
\noindent\textbf{FR-01 --- Course Administration.}
An Instructor must be able to create, update, and delete courses and organise their content into
named sections, while a Student must be able to browse and read any course they are enrolled in.

The Instructor creates a course with a title and description and controls whether it is published
and visible to students. Within a published course the Instructor defines \emph{sections}, each
with an ordered index that determines the display sequence. Course Materials (PDF documents,
slide sets, video files, or images) are uploaded into sections; the backend infers the file type
from the MIME type and extension and stores it alongside the file URL so the frontend can render
the appropriate viewer without an additional request. Enrollment is managed per student: the
Instructor can activate or drop individual students, and a student sees only the courses in which
their enrollment status is \texttt{ACTIVE}.

\bigskip
\noindent\textbf{FR-02 --- Content Extension System.}
A Student must be able to clone a course into a Private Workspace, add or modify materials
independently, and submit a Merge Proposal for the Instructor to review and optionally publish.

When a student clones a course, the backend performs a deep copy of the course structure into a
new \texttt{StudentWorkspace}: every section and every Course Material is mirrored as a
\texttt{WorkspaceSection} and a \texttt{WorkspaceMaterial}. Cloned materials carry the flag
\texttt{isReference = true}, meaning they point to the original file URL rather than a new copy,
which saves storage. The student may then upload additional materials into the Private Workspace.
When satisfied, the student submits a Merge Proposal, which the Instructor reviews in a dedicated
queue. On approval, the backend copies the student's file into the course storage and publishes it
as a new Course Material credited to the contributor.

\bigskip
\noindent\textbf{FR-03 --- Contextual Messaging System.}
Users must be able to post messages in the workspace chat and attach a direct reference to a
specific Course Material, so that other users can open the material from within the message.

The chat panel is embedded in the workspace view alongside the material list. Users compose
messages in a rich-text input and can type \texttt{@} followed by a material name to trigger an
autocomplete dropdown. Selecting a material embeds a non-editable tag chip in the message. When
the message is submitted, the chip is serialised into a portable token of the form
\texttt{@[id:title]}, which is stored as the message text. On render, the token is parsed back into
a clickable badge, which the user clicks to download or open the referenced Course Material. This
Contextual Anchor persists in the chat history and remains accessible to every enrolled student who
opens the workspace later.

\bigskip
\noindent\textbf{FR-04 --- AI-Based Content Querying.}
Users must be able to ask natural-language questions about the content of the currently open
workspace and receive answers grounded in that content.

The AI panel is a collapsible sidebar within the workspace view. When the user submits a question,
the frontend sends the question together with the URL of the selected Course Material to the
backend. The backend extracts the plain text from the document, constructs a structured prompt
that includes a system instruction, the document context truncated to 10\,000 characters, and the
user's question, then forwards the prompt to the Large Language Model (LLM) API. The response is
streamed back to the frontend and displayed with a typewriter animation. Conversation history is
preserved in the database through a hybrid memory model: a rolling buffer of recent messages and a
compressed long-term summary, ensuring that the AI can reference earlier turns in a session even
after the short-term buffer overflows.

\bigskip
\noindent\textbf{FR-05 --- Access Control.}
The system must recognise two roles, Student and Instructor, and enforce role-based restrictions
at the API level so that students cannot perform instructor-only operations.

Role assignment happens at registration and is stored as the \texttt{role} field of the
\texttt{User} entity. The authentication flow issues a signed JSON Web Token (JWT) upon successful
login; the token encodes the user's email and is valid for ten hours. Every subsequent request must
carry the token in the \texttt{Authorization: Bearer} header. The \texttt{JwtAuthenticationFilter}
intercepts each request, extracts and validates the token, and populates the Spring Security
context. Instructor-only endpoints are protected with \texttt{@PreAuthorize} annotations or
explicit ownership checks in the service layer so that a student account receives an HTTP 403
response if it attempts a restricted operation.

\subsection{Non-Functional Requirements}
\label{subsec:nonfunctional-requirements}

The following non-functional requirements define the quality constraints the system must satisfy
regardless of the specific feature being used.

\begin{table}[htbp]
    \centering
    \caption{Non-functional requirements}
    \label{tab:nonfunctional-requirements}
    \renewcommand{\arraystretch}{1.4}
    \begin{tabularx}{\textwidth}{|p{1.5cm}|X|}
        \hline
        \textbf{ID} & \textbf{Description} \\
        \hline
        NFR-01 & \textbf{Usability}: The interface should be straightforward enough for
                 students with no prior training to use without a manual. \\
        \hline
        NFR-02 & \textbf{Security}: All endpoints except login and registration must
                 require a valid JWT token. Role-based access must be enforced on the
                 server side. \\
        \hline
        NFR-03 & \textbf{Performance}: Course and workspace pages should load within
                 three seconds under normal conditions. \\
        \hline
        NFR-04 & \textbf{Maintainability}: The backend and frontend must be clearly
                 separated and independently deployable. The project structure should
                 follow established conventions for its respective framework. \\
        \hline
        NFR-05 & \textbf{Deployability}: The entire application must be runnable with
                 a single Docker Compose command without manual environment setup. \\
        \hline
    \end{tabularx}
\end{table}

% ----------------------------------------------------------------
\section{Use Cases}
\label{sec:use-cases}

The system has two types of actors: \emph{Student} and \emph{Instructor}. The use cases below
describe the main interactions each actor has with the system, grouped by role. Pre-conditions
are stated for the non-trivial cases that have dependencies on prior system state.

% TODO: Insert use case diagram here
% \begin{figure}[htbp]
%     \centering
%     \includegraphics[width=0.9\textwidth]{images/fig_2_3_use-case-diagram.png}
%     \caption{UML use case diagram for StudySpace}
%     \label{fig:use-case-diagram}
% \end{figure}

\subsection*{Instructor Use Cases}

\begin{itemize}
    \item \textbf{\texttt{Create course}}: The Instructor creates a course with a title and
          description and sets its published status.
    \item \textbf{\texttt{Upload Course Material}}: Pre-condition: the course and at least one
          section exist. The Instructor selects a file, assigns a title, and uploads it into a
          named section. The backend infers the file type and persists the file URL.
    \item \textbf{\texttt{Manage enrolments}}: The Instructor approves or drops individual
          student enrollments from the course management panel.
    \item \textbf{\texttt{Review Merge Proposal}}: Pre-condition: at least one student has
          submitted a Merge Proposal for this course. The Instructor inspects the proposed
          Course Material and either approves it (publishing it on the main course page with the
          contributor's name attached) or rejects it with an optional review note.
\end{itemize}

\subsection*{Student Use Cases}

\begin{itemize}
    \item \textbf{\texttt{Browse course materials}}: The student navigates to an enrolled course
          and opens individual sections to read the uploaded Course Materials.
    \item \textbf{\texttt{Clone course to Private Workspace}}: Pre-condition: the student is
          enrolled in a published course. The student triggers a clone action; the backend creates
          a Private Workspace containing mirror copies of all sections and Course Materials.
    \item \textbf{\texttt{Submit Merge Proposal}}: Pre-condition: the student owns a Private
          Workspace with at least one non-reference material or a modified reference material.
          The student selects materials and a target section, optionally adds a message, and
          submits the proposal for Instructor review.
    \item \textbf{\texttt{Post contextual message}}: The student types a message in the workspace
          chat, uses \texttt{@} to attach a Contextual Anchor to a specific Course Material, and
          submits. All enrolled users can see the message and follow the anchor.
    \item \textbf{\texttt{Query AI assistant}}: The student opens the AI panel, selects a Course
          Material for context, and types a question. The system sends the question and the
          extracted document text to the LLM API and displays the answer.
\end{itemize}

% ----------------------------------------------------------------
\section{Existing Solutions}
\label{sec:existing-solutions}

Several platforms address parts of the problem described in section~\ref{sec:problem-statement}.
This section compares the most relevant ones against the core features that StudySpace provides,
concluding with a synthesis of the integration gap that none of them fills end-to-end.

\begin{table}[htbp]
    \centering
    \small
    \caption{Comparison of existing platforms against StudySpace features}
    \label{tab:platform-comparison}
    \renewcommand{\arraystretch}{1.4}
    \begin{tabularx}{\textwidth}{>{\raggedright\arraybackslash}p{1.5cm} *{5}{>{\centering\arraybackslash\hspace{0pt}}X}}
        \toprule
        \textbf{Platform}
            & \textbf{Structured course materials}
            & \textbf{Student content contribution}
            & \textbf{Contextual discussion}
            & \textbf{AI content assistance}
            & \textbf{Role-based access control} \\
        \midrule
        \textbf{Zoom / Teams}
            & None & None & None & None & Partial \\
        \textbf{Discord}
            & None & None & None & None & Partial \\
        \textbf{Moodle}
            & Full & Limited & None & None & Full \\
        \textbf{Notion / Docs}
            & Full & Partial & None & Generic & None \\
        \textbf{StudySpace}
            & Full & Full & Full & Context & Full \\
        \bottomrule
    \end{tabularx}
\end{table}

The comparison in Table~\ref{tab:platform-comparison} shows that no existing platform covers all
areas simultaneously. The following subsections analyse each platform in detail.

% ----------------------------------------------------------------
\subsection{Zoom and Microsoft Teams}
\label{subsec:zoom-teams}

Zoom and Microsoft Teams are the dominant platforms for synchronous online education. Both support
high-quality video conferencing, screen sharing, breakout rooms, and session recording. Microsoft
Teams adds persistent channels, a file-sharing area backed by SharePoint, and deep integration
with the Microsoft 365 productivity suite.

Despite these strengths, neither platform is designed for asynchronous, material-centric
collaboration. Files shared during a meeting are stored in a flat SharePoint folder rather than
being organised into a course structure that reflects the academic curriculum. Discussion threads
are tied to meeting events, making conversations about a lecture slide inaccessible to students
who join the course later. There is no mechanism for students to annotate shared files and propose
changes back to the Instructor; the Instructor remains the sole publisher of content. The
asynchronous knowledge-management and student contribution use cases are left entirely unaddressed.

% ----------------------------------------------------------------
\subsection{Discord}
\label{subsec:discord}

Discord was originally designed for gaming communities but has been widely adopted by student
groups as an informal alternative to institutional platforms. Its server and channel model allows
students to self-organise into topic-specific channels, and its low barrier to entry means that
most students already have an account.

However, Discord's design is fundamentally chat-centric and lacks any concept of structured
educational content. Files are posted as raw message attachments in a general channel with no link
to a course section or curriculum structure. There is no role hierarchy that distinguishes
Instructors from students at the permission level beyond basic moderation roles; an Instructor
cannot publish a set of Course Materials in a read-only area that students can then clone and
annotate. Content disappears into scroll history, there is no access control aligned with course
enrolment, and there is no mechanism for structured content contribution or AI-assisted querying
of Course Materials.

% ----------------------------------------------------------------
\subsection{Moodle}
\label{subsec:moodle}

Moodle is a mature, open-source Learning Management System (LMS) used by thousands of
universities worldwide, including CTU. It provides the most comprehensive course management
feature set of any platform considered: Instructors can define courses with hierarchical sections,
upload files and links, set up graded assignments and quizzes, track student progress, and manage
enrolment through a role-based permission model.

Despite this breadth, Moodle has well-documented limitations with respect to student agency and
contextual collaboration. Its discussion model is based on forums that are standalone modules
attached to a course but not anchored to specific Course Materials; a student cannot attach a
comment directly to a particular document without navigating to a separate forum module. Student
contribution in Moodle is limited to wiki modules and assignment submissions, which are graded
artefacts that do not become part of the course content visible to other students. Moodle also
has no built-in AI assistant that operates on the actual course content in real time.

% ----------------------------------------------------------------
\subsection{Notion and Google Docs}
\label{subsec:notion-docs}

Notion and Google Docs represent a different category: general-purpose collaborative writing
tools. Both support real-time co-editing, rich text formatting, embedded media, and comment
threads. Notion extends this with a hierarchical page structure and customisable database views,
making it possible to approximate a course workspace.

The critical gaps are the absence of role-based access control aligned with academic roles and
the lack of a structured contribution workflow. In a shared Notion workspace, any collaborator
can edit any page, creating conflicts and the risk of accidental data loss in a course setting
with many students. Google Docs has stricter permission levels but no multi-tier workflow where a
student creates a private draft, develops it independently, and then submits it for review and
publication. Neither platform is AI-aware in a course-content sense: Notion AI and Google's
Gemini integration can generate or summarise text, but they operate on the open document in
isolation, without knowledge of which course the document belongs to or what other materials are
in the same section.

% ----------------------------------------------------------------
\subsection{Summary of Gaps}
\label{subsec:gap-summary}

The analysis above reveals a consistent pattern across all evaluated platforms. None of them
provides a mechanism for anchoring a discussion message to a named Course Material within a
structured course hierarchy, meaning that contextual references must be described manually in
text and the connection to the source material is easily lost. No platform supports a
clone-then-propose contribution model in which a student creates a private enriched copy of a
course section and submits it back for Instructor review and publication; the Instructor remains
the sole publisher of course content in every platform examined. Where AI assistance is
available, as in Notion AI or Google's Gemini integration, it operates on the open document in
isolation and cannot be grounded in the broader course context. The result is that satisfying all
four requirements simultaneously forces students to use at least two to three separate platforms,
reproducing the context-switching overhead described in section~\ref{sec:problem-statement}.

StudySpace addresses this integration gap by combining structured course management, a
clone-based student contribution workflow with Instructor review, chat messages with Contextual
Anchors tied to specific Course Materials, and an AI assistant grounded in the actual course
content, all within a single application with role-based access control.
